\subsection{Twisted \texorpdfstring{$A_5$}{A5} chief extensions over arbitrary finite bases}

\begin{statement}{D.30}{Proved}{Exact peripheral count invariance for every finite $A_5$ chief extension}
Let
\[
 1\longrightarrow A\longrightarrow E\xrightarrow{\pi}D
 \longrightarrow1,\qquad A\simeq A_5^{\Omega},
 \tag{D.30.1}
\]
be a finite extension in which $A$ is a nonabelian chief
factor.  Let $q_E:\widehat F\twoheadrightarrow E$ be an
epimorphism, put $X_i=q_E(x_i)$, $q_0=\pi q_E$, and
$d_i=q_0(x_i)$.  Let $q_1:\widehat F\twoheadrightarrow D$
be any epimorphism, and suppose its three specified
peripheral values are
\[
 q_1(x_i)=h_i d_i^\lambda h_i^{-1},\qquad
 h_i\in D,\quad \lambda\in\widehat\Z^\times.
 \tag{D.30.2}
\]
For arbitrary chosen lifts $H_i\in E$ of $h_i$, there
is an epimorphism $f:\widehat F\twoheadrightarrow E$
with
\[
 \pi f=q_1,\qquad
 f(x_i)=a_iH_iX_i^\lambda H_i^{-1}a_i^{-1}
                  \quad(a_i\in A).
 \tag{D.30.3}
\]
The number of these literal epimorphisms equals the
number for $(q_1,\lambda,h_1,h_2,h_3)=(q_0,1,1,1,1)$.
In particular it is positive.  There is no restriction
on the order of $D$, the degree of its factor action,
its fixed peripheral points, or its outer action on
$A$.  The result applies at every successive $A_5$
chief step of a finite tower with its given original
source epimorphism and its prescribed new peripheral
base epimorphism.
\end{statement}

\begin{proof}
Write $S=A_5$, $\widetilde S=S_5$, let
$\sigma=(01)\in\widetilde S$, and write
$\tau(s)=\sigma s\sigma^{-1}$.
The exact finite calculation in D.28 identifies
$\widetilde S\xrightarrow{\operatorname{Ad}}
\operatorname{Aut}(S)$ as an isomorphism and
$\operatorname{Out}(S)$ with $\mathbf F_2$ by parity.
In particular $C_{\widetilde S}(S)=1$.

\paragraph{The actual twisted permutation extension.}
The factor action of $D$ on $\Omega$ is transitive
because the kernel is chief.  Choose a factor
$\omega_0$, let $D_*$ be its stabilizer, and identify
$\Omega=D_*\backslash D$ with the right action.
Choose representatives $t_\omega\in D$ with
$t_{\omega_0}=1$.  The outer action of $D$ is a
homomorphism to
$\mathbf F_2^{\Omega}\rtimes\operatorname{Sym}(\Omega)$,
so its component $b$ satisfies
\[
 b_{de}(\omega)=b_d(\omega)+b_e(\omega d).
 \tag{D.30.4}
\]
Set
\[
 \epsilon(c)=b_c(\omega_0)\quad(c\in D_*),\qquad
 s_\omega=b_{t_\omega}(\omega_0),\qquad
 k(\omega,d)=t_\omega d t_{\omega d}^{-1}.
 \tag{D.30.5}
\]
Then $\epsilon:D_*\to\mathbf F_2$ is a character,
$k(\omega,d)\in D_*$, and the cocycle identity gives
\[
 b_d(\omega)=s_\omega+
        \epsilon(k(\omega,d))+s_{\omega d}.
 \tag{D.30.6}
\]
This formula supplies the required change of
coordinates, without averaging.  More explicitly,
centerlessness gives the isomorphism
\[
 E\longrightarrow
 \operatorname{Aut}(A)\times_{\operatorname{Out}(A)}D,
 \qquad e\longmapsto(\operatorname{Ad}_e,\pi(e)):
 \tag{D.30.7}
\]
injectivity and surjectivity follow exactly by the
kernel and inner-automorphism calculation in
\textup{(D.29.6)}.  Represent each factor automorphism
by its unique element of $\widetilde S$.  Conjugate
the resulting wreath coordinates by the vector
$(\sigma^{s_\omega})_\omega$.  In coordinates this
sends $(z,d)$ to
$(\sigma^{-s_\omega}z(\omega)\sigma^{s_{\omega d}},d)$.
Equations \textup{(D.30.6)--(D.30.7)} identify the
given extension, by this actual isomorphism, with
\[
 \begin{split}
 E_*={}&\{(z,d)\in\widetilde S^{\Omega}\rtimes D:
       \operatorname{sgn}z(\omega)
                   =\epsilon(k(\omega,d))\},\\
 (z,d)(w,e)={}&(z(\omega)w(\omega d),de).
 \end{split}
 \tag{D.30.8}
\]
Its kernel is precisely $S^{\Omega}$.  All the
$X_i,H_i$ are transported by the displayed
isomorphism; these transported elements are used
below.

For either source epimorphism $q=q_j$, put
$K_j=\ker q_j$ and $H_{j,C}=q_j^{-1}(C)$ for
$D_*\leq C\leq D$.  Choose actual words
$T_\omega\in\widehat F$ with $q(T_\omega)=t_\omega$
and $T_{\omega_0}=1$.  The twisted nonabelian
Shapiro map is the literal formula
\[
 \begin{split}
 \varphi:H_{j,D_*}\longrightarrow\widetilde S,\quad
 &\operatorname{sgn}\varphi=\epsilon q_j,\\
 z_\varphi(g)(\omega)
   & =\varphi(T_\omega gT_{\omega q_j(g)}^{-1}),\qquad
 f_\varphi(g)=(z_\varphi(g),q_j(g)).
 \end{split}
 \tag{D.30.9}
\]
The arguments of $\varphi$ lie in $H_{j,D_*}$;
multiplication and the parity condition in
\textup{(D.30.8)} follow by cancellation of the
middle representative.  Conversely, evaluate a
cocycle at $\omega_0$ on this subgroup.  Conjugating
by the unique vector $v\in S^{\Omega}$ with
$v(\omega_0)=1$ and
$v(\omega)=z(T_\omega)(\omega_0)$ makes it exactly
\textup{(D.30.9)}.  The entries of $v$ are even:
the transition element for $(\omega_0,t_\omega)$
in \textup{(D.30.5)} is $1$.
Thus there are exactly $60^{|\Omega|-1}$ literal
cocycles for each normalized surface map.

\paragraph{All intermediate extensions and their exact embeddings.}
For $D_*\leq C\leq D$, consider each character
$\epsilon_C:C\to\mathbf F_2$ extending $\epsilon$.
There may be none.  Choose representatives
$R_\nu$ for $\nu\in C\backslash D$, with the
identity at the base coset, and define
\[
 k_C(\nu,d)=R_\nu dR_{\nu d}^{-1},\qquad
 b_C(d,\nu)=\epsilon_C(k_C(\nu,d)).
 \tag{D.30.10}
\]
Let $E_C$ be the analogue of \textup{(D.30.8)}
with these data.  Write $p(\omega)=Ct_\omega$ and
\[
        g_\omega=\epsilon_C(t_\omega R_{p(\omega)}^{-1}).
 \tag{D.30.11}
\]
The argument belongs to $C$.  The following is an
injective homomorphism over $D$:
\[
 j_C:E_C\longrightarrow E_*,\qquad
 j_C(z,d)(\omega)=
       \sigma^{g_\omega}z(p(\omega))
                       \sigma^{-g_{\omega d}}.
 \tag{D.30.12}
\]
The base coordinate $d$ is retained.  Multiplication
is verified by cancelling the two adjacent powers
of $\sigma$; injectivity follows by reading one
coordinate above every $\nu$.  Its parity is
correct because
\[
 \epsilon(k(\omega,d))=
 g_\omega+b_C(d,p(\omega))+g_{\omega d}.
 \tag{D.30.13}
\]
This proves the embedding with its outer twisting.
At the base coordinate $g_{\omega_0}=0$; consequently
the Shapiro map of a homomorphism
$H_{j,C}\to\widetilde S$, followed by $j_C$, restricts
at $H_{j,D_*}$ to that very homomorphism's restriction.

\paragraph{The boundary constraint sets.}
Fix a peripheral $i$.  An element of $E_*$ over
$q_j(x_i)$ is conjugate under $A=S^{\Omega}$ to
another such element exactly when all their cycle
products are conjugate under $S$.  Indeed the
usual successive conjugator equation along a cycle
determines every conjugator after its first value;
all these values are even because the two component
parity patterns agree.  The last equation is exactly
the stated $S$-conjugacy of the cycle product.

This gives finite subsets $T_{j,C,\delta}$ of
$\widetilde S$, one for every boundary component
$\delta$ of $H_{j,C}$, that express the condition
on $j_Cf_\varphi(x_i)$ literally.  Here is an
explicit construction.  Let $\delta$ be a cycle
on $C\backslash D$, choose its initial coordinate
$\nu$, and choose the initial coordinate $\omega$
of each cycle $\gamma$ above it so that
$p(\omega)=\nu$.  If their lengths are $e$ and
$re$, respectively, the cycle product under
\textup{(D.30.12)} is
$\sigma^{g_\omega}a^r\sigma^{-g_\omega}$, where
$a$ is the product at $\delta$.  Intersect the
conditions
\[
 \sigma^{g_\omega}a^r\sigma^{-g_\omega}
       \in\mathcal C_{j,\gamma}
 \tag{D.30.14}
\]
over all such $\gamma$, and impose the parity
$\operatorname{sgn}a=\epsilon_C(q_j(c_\delta))$.
Here $\mathcal C_{0,\gamma}$ is the actual
$S$-class of the cycle product of $X_i$, and
$\mathcal C_{1,\gamma}$ comes from
$H_iX_i^\lambda H_i^{-1}$.
The result is $T_{j,C,\delta}$.  It is a union
of $S$-conjugacy classes, all of its prescribed
parity.  Formula \textup{(D.30.14)} is also valid
when the corresponding subset is empty.

The comparison of these sets retains the prescribed
$h_i$.  Match the boundary cycles by
\[
 C d\langle q_1(x_i)\rangle
       \longmapsto C d h_i\langle d_i\rangle.
 \tag{D.30.15}
\]
It preserves lengths and every relative boundary
degree, by \textup{(D.30.2)}.  Choose an initial
coordinate $\nu'$ of the new cycle and use
$\nu' h_i$ as the initial coordinate of its old
match.  Let $L_i\in E_C$ be the canonical element
over $h_i$, with components
$\sigma^{b_C(h_i,\nu)}$.
Its image under $j_C$ differs from $H_i$ by an
element of $A$.  Therefore conjugating a local
element by $L_i$ after taking its $\lambda$th
power transports precisely the old allowable
set to the new one.  Its cycle product is
conjugated by the actual component
$\sigma^{b_C(h_i,\nu')}$ of $L_i$.  Consequently
\[
 T_{1,C,\delta'}=
       \tau^{b_C(h_i,\nu')}
          \bigl(T_{0,C,\delta}^{\lambda}\bigr).
 \tag{D.30.16}
\]
The inverse power and inverse conjugation prove
the reverse inclusion, as $\lambda$ is a unit
on every finite group involved.  This argument
does not require $X_i$ itself to belong to
$j_C(E_C)$.

\paragraph{Actual surface maps with a specified parity character.}
The covers $H_{0,C}$ and $H_{1,C}$ have the same
index, genus, and boundary count, and
\textup{(D.30.15)} matches their boundary parities
for the characters $\chi_j=\epsilon_Cq_j$.
Both characters are trivial when $\epsilon_C=0$;
both are nontrivial otherwise, since $q_j$ maps
$H_{j,C}$ onto $C$.
There is an isomorphism of their surface groups,
and hence of their profinite completions, that
matches these characters and the indicated
oriented boundary classes.  The following gives
actual homeomorphisms realizing the comparison.
First use the adjacent Hurwitz maps
\textup{(D.29.13)} for the boundary permutation.
If all boundary character values vanish, the
remaining nonzero characters are covectors on
the handle symplectic space over $\mathbf F_2$.
A Dehn twist around a simple curve of class $v$
acts by the transvection $w\mapsto w+\langle w,v\rangle v$.
Every nonzero mod-two handle class has such a
curve: in each handle take the indicated one
of the three slopes $(1,0),(0,1),(1,1)$, and
join the selected disjoint curves by bands
away from the boundary.  These transvections
are transitive on nonzero vectors.  Explicitly,
for distinct $a,b$ with $\langle a,b\rangle=1$
use the transvection at $a+b$; when this pairing
is zero choose $c$ pairing to one with both
and use the two transvections through $c$.
This gives a product of boundary-fixed Dehn
twists matching the two nonzero covectors.

If a boundary class $d$ has character value
one, take two disjoint curves of classes
$a$ and $a+d$ bounding a pair of pants with
that boundary.  The bounding-pair map
$T_{a+d}T_a^{-1}$ fixes the boundary and acts
on homology by
\[
                 w\longmapsto w+\langle w,a\rangle d.
 \tag{D.30.17}
\]
Thus it adds the handle covector
$\langle -,a\rangle$ to the character.  Choosing
the handle basis curves supplies every required
translation.  For genus zero the boundary
values already determine the character.
Base these homeomorphisms outside their supports.
They induce actual inverse maps on the surface
groups and their profinite completions.  A
boundary generator may be conjugated by a
based path, whose parity is specified by
$\chi_j$; this can apply one further power of
$\tau$ to its target class.  This is the only
additional boundary change in this comparison.

We now prove the exact target-count fact that
accounts for all those possible changes.
For a fixed nontrivial character $\chi$ of a
surface group, the number of epimorphisms to
$\widetilde S$ with parity $\chi$ and boundary
values in specified $S$-invariant subsets is
unchanged by applying $\tau$ independently to
any of the boundary subsets.

Here is an integral finite proof.  In the group
algebra of $\widetilde S$, write $[T]=\sum_{t\in T}t$
and define
\[
 K_{ab}=\sum_{\substack{u,v\in\widetilde S\\
                 \operatorname{sgn}u=a,\operatorname{sgn}v=b}}
                       [u,v],
 \qquad a,b\in\mathbf F_2.
 \tag{D.30.18}
\]
Here each group element in the sum is its basis
element, so multiplicities are retained.
Let $\Delta=[5A]-[5B]$, using the two $S$-classes
of five-cycles.  Exact multiplication of the
120 permutations gives
\[
 \begin{gathered}
 K_{01}=K_{10}=K_{11}=L,\\
 L\Delta=0,\qquad L(K_{00}-L)=0,\\
 [O]\Delta=0,\qquad [O](K_{00}-L)=0
       \quad\text{for each odd }S\text{-class }O.
 \end{gathered}
 \tag{D.30.19}
\]
The coefficients of $K_{00}$ and $L$ on the
even classes, per element, are
\[
 \begin{array}{c|rrrrr}
 \text{class}&1&2^2&3&5A&5B\\\hline
 K_{00}&300&32&63&65&65\\
 L&180&72&63&45&45.
 \end{array}
 \tag{D.30.20}
\]
They vanish on odd permutations.  All eight
$S$-class sums in $\widetilde S$ commute.
The complete integer coefficient vectors and
all products in \textup{(D.30.19)} are checked
by the exhaustive permutation certificate
specified below.

For genus $g$ and boundaries $c_1,\ldots,c_b$,
the number of all homomorphisms with character
$\chi$ and boundary subsets $T_i$ is exactly
the coefficient of $1$ in
\[
 \prod_{j=1}^g K_{\chi(a_j),\chi(b_j)}
                   \prod_{i=1}^b[T_i].
 \tag{D.30.21}
\]
This is the tuple count for the literal surface
relation $\prod[a_j,b_j]\prod c_i=1$.
Changing one even boundary subset by $\tau$
changes its class sum by $0$, $\Delta$, or
$-\Delta$; odd classes are unchanged.
Nontriviality of $\chi$ provides either a
factor $L$ or another odd boundary factor.
Thus \textup{(D.30.19)} makes the difference
in \textup{(D.30.21)} zero.

The same conclusion holds after excluding
nonsurjective maps.  The exhaustive subgroup
lattice of $S_5$ has 156 subgroups.  Exactly
six proper subgroups contain both an odd
permutation and a five-cycle, and each is
the normalizer $F_{20}=C_5\rtimes C_4$ of
its unique subgroup of order five.
In the difference from one changed boundary,
every contributing nonsurjective map has
precisely one of these six subgroups as its
image.  It has a five-cycle because the changed
factor is supported there, and an odd image
because $\chi$ is nontrivial.
For each embedded $F_{20}$ put $R=\sum_{c\in C_5}c$.
Its corresponding restricted distributions are
\[
 K^{F_{20}}_{00}=25\cdot1+15R,
 \qquad
 K^{F_{20}}_{01}=K^{F_{20}}_{10}
                   =K^{F_{20}}_{11}=20R.
 \tag{D.30.22}
\]
The restriction of $\Delta$ to $C_5$ has
augmentation zero.  Every nonempty odd
intersection of an $S$-class with $F_{20}$
is the union of its two odd $C_5$-cosets.
Consequently both $20R$ and every such odd
class sum annihilate that restricted
$\Delta$.  The relevant restricted class
sums commute; this also follows immediately
by multiplying in $C_5\rtimes C_4$ and is
checked for all six embeddings.
The difference of the homomorphism counts
into each $F_{20}$ is therefore zero by the
same formula \textup{(D.30.21)}.
Every map in this difference is already
onto that $F_{20}$, by the proved proper
subgroup classification.  There is no
overcount between the six image groups.
Subtracting their zero differences proves
the asserted equality of epimorphism counts
to $S_5$.

Finally every unit power on the eight
$S$-classes in $S_5$ is realized by the
actual automorphism $\tau_\lambda$: it is
the identity for $\lambda\equiv1,4\pmod5$
and $\tau$ for $\lambda\equiv2,3\pmod5$.
The three odd classes of orders $2,4,6$
are each preserved by all unit powers;
the even classes have the transformation
in \textup{(D.29.12)}.
Thus the surface isomorphism above,
$\tau_\lambda$, and the independent-flip
count prove equality of the following
two finite numbers, for each $C,\epsilon_C$:
\[
 P_{j,C,\epsilon_C}
 =\#\{\varphi:H_{j,C}\twoheadrightarrow G_C:
       \operatorname{sgn}\varphi=\epsilon_Cq_j,
       \ \varphi(c_\delta)\in T_{j,C,\delta}\},
 \quad
 G_C=\begin{cases}A_5&\epsilon_C=0,\\S_5&\epsilon_C\ne0.
       \end{cases}
 \tag{D.30.23}
\]
For $\epsilon_C=0$, the character and all
flips in \textup{(D.30.16)} and in the
based surface comparison are trivial;
the actual target automorphism and source
isomorphism give the equality directly.
For $\epsilon_C\ne0$, the just-proved
fixed-character count applies.  No claim
that the independent class flips themselves
are group automorphisms is needed.

\paragraph{The exact finite-base subtraction.}
An epimorphism counted in
\textup{(D.30.23)} has
$\varphi(K_j)=1$ or $A_5$.
Indeed this image is contained in $A_5$
by the parity condition and is normal in
$G_C$, hence normal in the simple group
$A_5$.  The first alternative consists
precisely of maps
\[
 \varphi=\Phi q_j|_{H_{j,C}},\qquad
 \Phi:C\twoheadrightarrow G_C,
              \operatorname{sgn}\Phi=\epsilon_C.
 \tag{D.30.24}
\]
Their admissible counts on the two sides
agree by an actual finite-group comparison.
For such a $\Phi$, form the induced section
\[
 s_\Phi:D\longrightarrow E_C,
 \qquad
 s_\Phi(d)(\nu)=\Phi(R_\nu dR_{\nu d}^{-1}),
 \tag{D.30.25}
\]
with base coordinate $d$.
It is a homomorphism by cancellation of
representatives, and its parity is
\textup{(D.30.10)}.  The two maps in
\textup{(D.30.24)} give exactly the
Shapiro maps $s_\Phi q_0$ and $s_\Phi q_1$.
Their compositions with $j_C$ have the
literal equality
\[
 j_Cs_\Phi(q_1(x_i))
  =j_Cs_\Phi(h_i)\,
       j_Cs_\Phi(d_i)^\lambda\,
                      j_Cs_\Phi(h_i)^{-1}.
 \tag{D.30.26}
\]
The element $j_Cs_\Phi(h_i)$ and the
prescribed $H_i$ have the same base
coordinate; their difference is in $A$.
It follows in both directions that the
same $\Phi$ is admissible for $j=0$ and
$j=1$.  The inverse unit power proves
the reverse implication.  This is a
bijection of the maps that factor through
the finite base; it also retains their
surjectivity and prescribed character.

Subtract this common number from
\textup{(D.30.23)}.  We obtain equal
numbers $N_{j,C,\epsilon_C}$ of maps
with the same boundary conditions and
$\varphi(K_j)=A_5$.  Put
\[
       N_{j,C}=\sum_{\epsilon_C|_{D_*}=\epsilon}
                             N_{j,C,\epsilon_C}.
 \tag{D.30.27}
\]
The empty sum is zero.  We have proved
$N_{0,C}=N_{1,C}$ for every $C$, without
assuming that a finite-base quotient of
$A_5$ is impossible.

\paragraph{Normalizer, generation, and literal counts.}
Fix $j$ and a map
$\varphi:H_{j,D_*}\to S_5$ of prescribed
character and boundary classes with
$\varphi(K_j)=A_5$.  Let
$N=\ker(\varphi|_{K_j})$.
Its normalizer in $\widehat F$ contains
$H_{j,D_*}$ and $K_j$, so it equals
$H_{j,B}$ for a unique $B\geq D_*$.
For $C\geq D_*$, the condition that
$H_{j,C}$ normalize $N$ is equivalent
to extension of $\varphi$ to a map
$H_{j,C}\to S_5$ with image containing
$A_5$.  The extension is unique.
For existence, conjugation on
$K_j/N\xrightarrow[\varphi]{\sim}A_5$
gives the actual homomorphism
$H_{j,C}\to\operatorname{Aut}(A_5)=S_5$.
Its restriction agrees with $\varphi$
because $C_{S_5}(A_5)=1$.
Its parity factors through a unique
character $\epsilon_C:C\to C_2$
extending $\epsilon$.  Conversely every
extension with $K_j$-image $A_5$
normalizes $N$, and centerlessness gives
uniqueness.  Thus each map whose
normalizer contains $H_{j,C}$ is counted
exactly once in \textup{(D.30.27)}.

Let $e_{j,B}$ count those whose normalizer
is exactly $H_{j,B}$.  On the same finite
subgroup poset on both sides,
\[
 N_{j,C}=\sum_{B\geq C}e_{j,B},\qquad
 e_{j,D_*}=\sum_{C\geq D_*}\mu(D_*,C)N_{j,C}.
 \tag{D.30.28}
\]
The already proved equalities of $N_{j,C}$
therefore give $e_{0,D_*}=e_{1,D_*}$.
Under \textup{(D.30.9)}, the coordinate
maps on $K_j$ are
$k\mapsto\varphi(T_\omega kT_\omega^{-1})$.
They are onto $A_5$, and their kernels
are pairwise distinct exactly when the
normalizer is $H_{j,D_*}$.
The simple-factor subdirect argument of
D.25 then gives the whole product
$A_5^{\Omega}$.  Together with the
epimorphism to $D$ this is precisely
surjectivity onto $E_*$.
The original epimorphism $q_E$ shows
$e_{0,D_*}>0$.
Finally the unique even gauges in
\textup{(D.30.9)} show that the literal
proper-lift count on side $j$ is
\[
            60^{|\Omega|-1}
              \sum_{C\geq D_*}\mu(D_*,C)N_{j,C}.
 \tag{D.30.29}
\]
The gauges preserve the required
$A$-conjugacy classes, so no further
quotient or centralizer factor is
implicit in this formula.
Its equality and positivity prove
\textup{(D.30.3)}.

\paragraph{Finite certificates and the scope of iteration.}
The script \texttt{a\_peripheral\_s5\_twist.py}
constructs all 120 permutations, enumerates
all 156 subgroups, and verifies every
integer coefficient in
\textup{(D.30.19)--(D.30.22)}.
It verifies all sixteen units modulo
$60$ on the eight $A_5$-classes in
$S_5$.  Its output and JSON certificate
end with
\texttt{PASS\_TWISTED\_S5\_CLASS\_ALGEBRA}.
The independent map script
\texttt{a\_peripheral\_s5\_maps.py}
uses the actual bases $D=S_3,S_4$,
the stabilizer generated by a transposition,
and every extending character on every
intermediate subgroup.  It verifies
4680 literal multiplication comparisons
for \textup{(D.30.12)} and 1473 nonempty
boundary comparisons for
\textup{(D.30.16)}, with all permutation
coordinates retained.  Its output and
JSON end with
\texttt{PASS\_TWISTED\_WREATH\_EXACT\_MAPS}.
The general maps and count argument above
prove the statement beyond those finite
test bases.

At a successive $A_5$ chief step, use the
epimorphism already obtained at the
previous base as $q_1$; its peripheral
values have exactly the required form.
The given original source epimorphism
at that step supplies $q_E$.  This
proves the asserted finite iteration.
It makes no assertion that towers with
these chief factors are cofinal among
all finite quotients, and it does not
remove the central obstruction in D.24
or other chief factors in D.1--D.2.
\end{proof}
